\documentclass[12pt]{article}

\usepackage[margin=1in]{geometry}
\usepackage{amsmath,amssymb}
\usepackage{enumitem}
\usepackage{booktabs}
\usepackage{longtable}
\usepackage[authoryear,round]{natbib}
\usepackage{hyperref}
\hypersetup{
    colorlinks=true,
    linkcolor=black,
    citecolor=black,
    urlcolor=black
}

\title{\textbf{Literature Review, Paper Summaries, and Search Strategies} \\ For: Endogenous Education Choice in Segmented Search Markets}
\author{Ramzi Chariag}
\date{March 2026}

\begin{document}
\maketitle
\tableofcontents
\newpage

%% ====================================================================
\part{Paper Summaries}
%% ====================================================================

%% ------------------------------------------------------------------
\section{\citet{nekoei2017} --- ``Does Extending Unemployment Benefits Improve Job Quality?''}
%% ------------------------------------------------------------------

\subsection*{Main point}

More generous unemployment benefits improve the quality of post-unemployment matches. The standard view holds that longer UI raises moral hazard (workers search less), but the paper shows an offsetting ``selectivity'' channel: workers with more insurance can afford to be choosier and wait for better offers.

\subsection*{Specification}

The authors exploit a sharp regression-discontinuity design in Austria. Workers who become unemployed after age~40 are entitled to 39~weeks of UI benefits; those just under~40 get only 30~weeks. This 9-week jump at an age threshold creates quasi-experimental variation in UI generosity that is unrelated to worker ability or match quality.

\subsection*{Key findings}

\begin{itemize}
    \item A 9-week UI extension increases reemployment wages by approximately 0.5\% on average (about 1.3\% per additional week in the more generous interpretation).
    \item The wage gain comes from improved match quality --- workers find jobs at better-fitting firms --- not from higher bargaining power within the same type of match.
    \item The wage effect is entirely concentrated among workers who find jobs \emph{after} the benefit extension kicks in (late exiters), not those who would have found jobs anyway (early exiters).
    \item There is a negative cross-study correlation between the unemployment duration effect and the wage effect of UI extensions: in labour markets where UI extends search duration substantially (weak labour markets), the wage gains are smaller, consistent with a duration-dependence force working against the selectivity force.
\end{itemize}

\subsection*{Mechanism}

Two offsetting forces: (i)~\emph{selectivity} --- workers raise their reservation match quality because the outside option during search is higher, leading to better matches; (ii)~\emph{duration dependence} --- longer search depletes savings, erodes skills, or signals lower quality to employers, pushing workers into worse matches. The net effect on wages depends on which force dominates. In the Austrian RD, the selectivity effect dominates.

\subsection*{Takeaways for your project}

\begin{enumerate}
    \item \textbf{Structural microfoundation.} Your model provides exactly the mechanism \citet{nekoei2017} document empirically. In your framework, raising $b_U$ increases the reservation productivity threshold $p^*(x,z)$, making unskilled workers more selective. This is the selectivity channel. Your model can quantify it structurally and decompose it into duration vs.\ quality components.

    \item \textbf{Cross-market feedbacks.} \citet{nekoei2017} work in a single-market setting. Your model adds the dimension they cannot capture: raising $b_U$ not only improves match quality in the unskilled market, but also shifts the training threshold $\bar{x}$ and changes the composition of the skilled unemployment pool, with further tightness and match-quality effects in the skilled market.

    \item \textbf{Outside option $\neq$ UI.} The paper uses UI extension as an instrument, but the underlying economics is about outside options. Your Policy~Exercise~(1) studies outside options generically ($b_U$, $b_S$) --- the \citet{nekoei2017} finding motivates the exercise without requiring you to label $b$ as UI.

    \item \textbf{The duration--quality trade-off.} Your analytical derivations of expected duration $D_j(x)$ and expected match quality $\mathbb{E}[p \mid p \geq p^*(x,z)]$ as functions of $b$ are precisely the objects needed to replicate the \citet{nekoei2017} decomposition structurally.
\end{enumerate}


%% ------------------------------------------------------------------
\section{\citet{barany2016} --- ``The Minimum Wage and Inequality: The Effects of Education and Technology''}
%% ------------------------------------------------------------------

\subsection*{Main point}

The minimum wage affects inequality not only through direct wage compression, but also through its effect on the \emph{education margin}: by changing skill prices, it alters the incentive to acquire education, which changes the skill composition of the workforce, which in turn changes the direction of technical change. The paper builds a GE model that captures all three channels simultaneously.

\subsection*{Specification}

Two-sector model (high-skilled, low-skilled) with:
\begin{itemize}
    \item Heterogeneous workers who differ in ability and time cost of education.
    \item Endogenous education choice: workers above an ability threshold choose to become high-skilled.
    \item Endogenous directed technical change: firms direct R\&D toward the skill type that is relatively more abundant/profitable, following Acemoglu (2002).
    \item A binding minimum wage in the low-skilled sector.
    \item Competitive labour markets (no search frictions).
\end{itemize}

The education decision is a threshold rule: workers with ability above $\bar{x}$ invest in education. The threshold is determined by the indifference condition where the return to education (higher wage in the skilled sector minus cost) equals zero at the margin.

\subsection*{How she motivates the education margin}

\begin{itemize}
    \item The paper argues that you cannot evaluate minimum wage effects without accounting for the education response. A higher minimum wage compresses the low-skilled wage distribution, which \emph{reduces} the return to education for some marginal workers (the low-skilled sector just got better), but \emph{also} changes the relative supply of skills, which affects directed technical change and thus the skill premium.
    \item The motivation is that the policy $\to$ education margin $\to$ composition $\to$ GE feedback chain is quantitatively important: ignoring the education response misses a significant part of the effect.
    \item She estimates the model and finds that the decline in the real minimum wage since the 1980s accounts for a meaningful share of rising wage inequality, once the education and technology channels are accounted for.
\end{itemize}

\subsection*{Relevance to your project}

\begin{enumerate}
    \item \textbf{Same spirit, different frictions.} \citet{barany2016}'s model has endogenous education choice with endogenous market composition and GE feedbacks --- the same conceptual architecture as your model. The key difference: she has competitive markets, you have search frictions. Your paper can position itself as bringing this education-margin logic into a frictional setting where the continuation values that govern the threshold are fundamentally richer.

    \item \textbf{Advisor connection.} Citing \citet{barany2016} is natural and expected --- it is your advisor's work and it has the closest conceptual structure. The literature review should note the paper's contribution and distinguish yours by the frictional environment.

    \item \textbf{Motivation approach.} Observe how she motivates: (a)~state the policy question, (b)~argue that the education margin matters for that question, (c)~show that existing models miss the feedback, (d)~build the minimal model that captures it. Your motivation follows the same template but replaces ``minimum wage'' with ``outside options / education subsidies'' and ``competitive markets'' with ``frictional markets.''
\end{enumerate}


%% ------------------------------------------------------------------
\newpage
\section{\citet{charlot2005} --- ``Self-Selection in Education with Matching Frictions''}
%% ------------------------------------------------------------------

\subsection*{Main point}

In a frictional two-sector labour market where workers self-select into education, the decentralised equilibrium \emph{always} produces over-education relative to the social optimum. The key mechanism is a composition externality: the marginal entrant into education does not internalise the deterioration she imposes on the mean ability of both the educated and uneducated pools, reducing job creation in both markets.

\subsection*{Model structure}

Steady-state, two-sector matching model. Workers are heterogeneous in a single innate ability $a \in [0,1]$ with density $\phi$. Two sectors (high $h$, low $l$) each have their own CRS matching function; tightness $\theta_i = v_i/u_i$ is sector-specific. Output in sector $i$ is $y_i(a) = A_i \cdot a$, with $A_h > A_l > 0$. Wages are set by Nash bargaining with bargaining power $\beta$; the resulting wage schedule is linear in ability. Free entry drives the vacancy value to zero, pinning each sector's tightness as a function of the mean ability of workers in that pool. The two markets are in \emph{simultaneous general equilibrium}: both free-entry conditions and the self-selection condition hold jointly. Job destruction arises only through an exogenous death/replacement rate $\delta$; there is no endogenous separation. Workers make a binary education choice at birth: pay a fixed, ability-independent cost $C$ (net of any subsidy $S$) to enter market $h$, or stay in market $l$ for free.

The equilibrium is characterised by a threshold $\sigma^* \in (0,1)$ such that workers with $a \geq \sigma^*$ educate. Because both tightness levels depend on the mean ability of each pool, which itself depends on $\sigma^*$, the equilibrium is a fixed point in threshold space.

\subsection*{Key findings}

\begin{enumerate}
    \item The unique efficient outcome features a \emph{tax} on education, not a subsidy. Over-education holds for all parametrisations (bargaining power, schooling cost, technology).
    \item The composition effect --- mean ability of both pools is endogenous to $\sigma^*$ and feeds back into job creation --- is the novel externality specific to frictional markets. It is absent in frictionless competitive settings.
    \item A rise in the schooling subsidy lowers $\sigma^*$, reducing mean ability in both pools, causing tightness and job-finding rates to fall in both sectors simultaneously.
    \item Multiple equilibria are possible when the multiplier (education lowers unskilled pool quality, cheapening the unskilled outside option, encouraging more education) dominates the stabiliser.
\end{enumerate}

\subsection*{Takeaways for your project}

\begin{enumerate}
    \item \textbf{Closest structural predecessor.} This paper is the most direct antecedent of your model's architecture: binary education choice, two markets in simultaneous frictional equilibrium, endogenous composition of both pools, and Nash bargaining throughout. Your model inherits this framework.

    \item \textbf{Key gap: no job ladder.} Charlot--Decreuse (2005) has no on-the-job search and no job ladder in the skilled market. Wages are determined at matching and never renegotiated due to outside competition. Your model adds a full job ladder (CPR 2006 style) in the skilled market, making the skilled sector's wage dynamics fundamentally richer and giving it a non-trivial internal wage distribution.

    \item \textbf{Key gap: no endogenous job destruction.} Separation occurs only through death, ruling out the reservation-quality mechanism that governs your unskilled market ($p_U^*$). Endogenous job destruction in your unskilled market is what makes the outside-option exercise ($b_U$) tractable: higher $b_U$ shifts $p_U^*$, not merely search intensity.

    \item \textbf{Key gap: ability-independent cost.} The cost $C$ is the same for all workers. Your training cost $c(x) = c_0(1-x)e^{-x}$ declines in ability, making the composition externality richer and the policy analysis of cost subsidies more nuanced.

    \item \textbf{Key gap: no aggregate states.} Your model adds transitions between aggregate states $z \in \mathcal{Z}$, so the training threshold $\bar{x}(z)$ is endogenous to the business cycle --- a dimension entirely absent here.

    \item \textbf{Correct the literature table.} The existing table states the gap as ``composition exogenous, markets not jointly in equilibrium,'' which is incorrect for this paper. Both the composition \emph{and} the simultaneous equilibrium are present here. The genuine gap is: no job ladder, no aggregate states, and an ability-independent schooling cost.
\end{enumerate}


%% ------------------------------------------------------------------
\section{\citet{charlot2005b} --- ``Adaptability, Productivity, and Educational Incentives in a Matching Model''}
%% ------------------------------------------------------------------

\subsection*{Main point}

When education raises both a worker's productivity within a job and her \emph{adaptability} (the fraction of technologies she can operate), the direction of the education-market inefficiency depends on which dimension dominates. Productivity-only education is under-supplied; adaptability-only education is over-supplied. In the general case, either can hold.

\subsection*{Model structure}

Continuous time, a continuum of technologies $j \in [0,1]$ each with its own sub-market but with identical tightness $\theta$ in symmetric equilibrium. Workers are \emph{homogeneous in ability} and choose a continuous schooling \emph{duration} $T_i$ rather than a binary type. Education raises: (i)~\emph{productivity} $y(T_i)$ on any matched job, and (ii)~\emph{adaptability} $F(T_i)$, the share of technologies the worker can operate, so her contact rate is $F(T_i) \cdot m(\theta)$. Education is purely time-consuming (opportunity cost only; no tuition). All workers choose the same $T^*$ in equilibrium. Wages are Nash bargaining with parameter $\beta$. Free entry pins market tightness $\theta^*$. The equilibrium is a pair $(T^*, \theta^*)$ satisfying simultaneously the Optimal Schooling condition and the Free-Entry condition. A unique equilibrium exists.

\subsection*{Key findings}

\begin{enumerate}
    \item When education raises only productivity: under-investment, because higher productivity raises vacancy creation (a positive externality workers do not internalise).
    \item When education raises only adaptability: over-investment, because the extensive wage return to adaptability (improving one's outside option in bargaining) is privately valuable but not socially so, and congestion in the additional sub-markets harms others.
    \item Both effects combined admit either sign for the inefficiency.
    \item Higher job destruction $q$ raises schooling duration when adaptability dominates --- consistent with European data from the 1970s--1980s showing simultaneous rises in education and unemployment.
\end{enumerate}

\subsection*{Takeaways for your project}

\begin{enumerate}
    \item \textbf{No pool segmentation.} All workers are identical; there is a single representative-agent optimum, not a threshold splitting a heterogeneous population into two markets. Your model's composition externality (endogenous mean ability of two pools feeding back into job creation in both) is absent here.

    \item \textbf{Continuous vs.\ binary training.} Your model features a binary threshold $\bar{x}(z)$, closer to Charlot--Decreuse (2005). The continuous duration choice here may be relevant if you introduce a schooling-length margin in future extensions, but is not part of your current structure.

    \item \textbf{No two-market simultaneous equilibrium.} While the market is segmented by technology, all workers compete in a single aggregate market with a single tightness. The two-market structure with separate equilibria is your paper's key contribution.

    \item \textbf{No job ladder.} No PVR sequential auction, no on-the-job search, no wage dynamics within a sector.

    \item \textbf{No aggregate states.}
\end{enumerate}


%% ------------------------------------------------------------------
\section{\citet{wasmer2006} --- ``General versus Specific Skills in Labor Markets with Search Frictions and Firing Costs''}
%% ------------------------------------------------------------------

\subsection*{Main point}

In a single frictional labour market with firing costs, the equilibrium skill regime (general or specific human capital) is determined by the level of search frictions: high frictions favour specific skills (which raise within-firm productivity at the cost of portability), while low frictions favour general skills (which preserve outside options at the cost of within-firm productivity). Employment protection (firing taxes) pushes economies toward specific-skill equilibria, rationalising transatlantic differences in turnover and training.

\subsection*{Model structure}

Standard Mortensen--Pissarides (1994) extension to a \emph{single} labour market. Workers are \emph{homogeneous in ability}; each newborn chooses between specific ($s$) or general ($g$) human capital at job entry. Production: $y = \varepsilon + h^k$, where $\varepsilon$ is an idiosyncratic match-specific productivity shock (Poisson arrival rate $\lambda$, drawn from distribution $F$) and $h^k$ is the worker's human capital level ($h^s > h^g$ within firm; general capital is retained across jobs, specific capital is lost on separation). Wages are Nash bargaining at two stages: an entry wage and a renegotiated continuation wage after shocks. Job destruction is endogenous: a match is destroyed when surplus $S_k(\varepsilon) < 0$, defining a reservation productivity $R^k$. Free entry pins tightness $\theta^K$. Three regimes are possible: $S$ (all specific), $G$ (all general), and $M$ (mixed). A firing tax $T$ (paid by the firm, non-transferable to worker) is included.

\subsection*{Key findings}

\begin{enumerate}
    \item Workers prefer specific skills when job-finding rate $p$ is low and general skills when $p$ is high. The threshold $p^M$ separates the two pure regimes.
    \item Specific skills and low turnover reinforce each other (complementarity): the S-regime features longer job duration and higher within-firm productivity.
    \item Higher employment protection (firing tax $T$) raises the likelihood of the S-regime by reducing tightness $\theta$ and hence lowering $p$.
    \item Application: continental Europe's higher firing costs $\to$ specific skills $\to$ lower turnover $\to$ higher long-term unemployment.
\end{enumerate}

\subsection*{Takeaways for your project}

\begin{enumerate}
    \item \textbf{Single market, not two markets.} This is the key distinction from your paper. Wasmer (2006) has one labour market in which all workers --- regardless of skill type --- compete. There is no separation into a skilled and an unskilled market with different tightness levels, different wage distributions, and separate free-entry conditions. The existing literature table's description (``two markets not simultaneously in equilibrium'') is misleading: the correct characterisation is that this paper has \emph{no market segmentation at all}.

    \item \textbf{No composition externality.} The fraction $\kappa^g$ in the mixed regime is determined by a representative-agent indifference condition, not by a threshold over a heterogeneous distribution. There are no pools whose mean characteristics feed back into job creation.

    \item \textbf{No job ladder.} Workers stay in their current firm until idiosyncratic shocks trigger separation. There is no PVR sequential auction, no on-the-job wage dynamics, no job-to-job mobility. Wage renegotiations respond only to idiosyncratic $\varepsilon$ shocks, not to outside offers.

    \item \textbf{On-the-job vs.\ pre-market training.} The baseline model has workers choosing skills at job entry. Although Section~IVB shows the schooling extension produces identical steady states, the fundamental structure remains a single-market model.

    \item \textbf{Endogenous job destruction.} Wasmer (2006) is a useful reference for the endogenous job-destruction mechanism in your \emph{unskilled} market, where the DMP reservation-quality rule $p_U^*$ plays an analogous role to $R^k$ here.

    \item \textbf{No aggregate state transitions.}
\end{enumerate}


%% ------------------------------------------------------------------
\section{\citet{charlot2010} --- ``Over-Education for the Rich, Under-Education for the Poor''}
%% ------------------------------------------------------------------

\subsection*{Main point}

When workers are heterogeneous in \emph{both} ability and schooling cost, the decentralised frictional equilibrium simultaneously over-educates low-ability, low-cost (``rich'') workers and under-educates high-ability, high-cost (``poor'') workers. The mechanism is a market adverse selection: rich workers crowd the educated pool, reducing its mean ability and depressing job creation in the high-skill sector, which undermines the return to schooling for the talented poor.

\subsection*{Model structure}

Direct extension of \citet{charlot2005} with \emph{two-dimensional worker heterogeneity}: ability $a$ and schooling cost $c$, jointly distributed with CDF $\Phi$ on $\mathbb{R}_+^2$. Two sectors ($h$, $l$) with $y_{ia} = A_i \cdot a$, $A_h > A_l$; each has a CRS matching technology with tightness $\theta_i$. Search in each sector is random (ability is unknown to firms ex ante), so the free-entry condition in sector $i$ depends on mean ability $\bar{a}_i$ of the job-seekers in that pool. Wages are Nash bargaining; job destruction is exogenous (death rate $\delta$). The two markets are in simultaneous general equilibrium. The education decision: worker $(a,c)$ chooses education if $U_{ha} - c \geq \max\{U_{la}, 0\}$. This defines a linear indifference locus $c = R^* \cdot a$ in the $(a,c)$ plane, where the endogenous return $R^*$ depends on tightness in both sectors. The equilibrium is a fixed point in mean abilities $(\bar{a}_h, \bar{a}_l)$.

\subsection*{Key findings}

\begin{enumerate}
    \item Simultaneous over-education for the rich ($a$ low, $c$ low) and under-education for the poor ($a$ high, $c$ high), provided ability and schooling cost are not too positively correlated.
    \item The private return $R^*$ is independent of ability (a flat indifference line in $(a,c)$ space), while the social return $R^s(a)$ is increasing in $a$. Hence the rich over-invest and the poor under-invest.
    \item Optimal policy is redistributive: a subsidy proportional to schooling cost, combined with a lump-sum tuition fee. Net, education should be taxed, not subsidised.
    \item Generalises \citet{charlot2005}: with degenerate cost distribution, the model collapses to universal over-education.
\end{enumerate}

\subsection*{Takeaways for your project}

\begin{enumerate}
    \item \textbf{Two-dimensional heterogeneity.} Your model features ability heterogeneity (type $x$) and an ability-dependent training cost $c(x) = c_0(1-x)e^{-x}$. This creates a single-dimensional threshold $\bar{x}(z)$: unlike the 2010 paper's linear locus in $(a,c)$ space, your threshold is a scalar. The key distinction is that your cost is a \emph{known function} of type, not an independent draw, so there is no two-dimensional selection problem.

    \item \textbf{No job ladder in the skilled market.} As in Charlot--Decreuse (2005), the high-skill sector here operates with static Nash bargaining and no PVR wage dynamics. Your model's job ladder is the central innovation not present in this paper.

    \item \textbf{Simultaneous two-market equilibrium.} Your paper shares this with Charlot--Decreuse (2010). The key addition is the job ladder and aggregate states.

    \item \textbf{Policy exercise connection.} The 2010 paper's analysis of redistributive education subsidies --- differentiating between ability and cost dimensions --- is directly relevant to your Policy~Exercise~(2): comparing flow training subsidies (reducing $b_T$) vs.\ upfront cost reductions (reducing $c_0$). Your framework allows you to quantify both effects in a calibrated model.

    \item \textbf{No aggregate states.} Charlot--Decreuse (2010) is a steady-state model. Your model's aggregate states are the source of the cyclical training threshold, absent here.
\end{enumerate}


%% ------------------------------------------------------------------
\section{Charlot and Malherbet (2013) --- File Mismatch Note}
%% ------------------------------------------------------------------

\subsection*{Note on the uploaded file}

The file \texttt{Charlot-Malherbet2013.pdf} in the workspace does \textbf{not} contain the intended paper. Instead, the file is:

\begin{quote}
Charlot, Olivier, Franck Malherbet, and Mustafa Ulus. 2013. ``Efficiency in a Search and Matching Economy with a Competitive Informal Sector.'' \emph{Economics Letters} 118(1): 192--194.
\end{quote}

This is a 3-page note on informality (comparing formal search to informal employment as an outside option) and has no education or training content. It is largely orthogonal to the Education-in-Search family and shares only the general Pissarides framework.

\medskip

The intended paper --- Charlot, Olivier, and Franck Malherbet (2013). ``Education and Self-Selection in a Search Economy.'' \emph{Research in Economics} --- has \textbf{not yet been read} and a summary cannot be provided. The correct file needs to be obtained (e.g.\ via ScienceDirect, author websites, or REPEC). Once the correct paper is in hand, a summary following the standard format should be added here.

\medskip

\noindent\textbf{Action required:} Locate and download the correct Charlot--Malherbet (2013) \emph{Research in Economics} paper and add its summary before circulating this document.


%% ====================================================================
%% FAMILY 3: SEARCH FOUNDATIONS
%% ====================================================================

%% ------------------------------------------------------------------
\section{\citet{diamond1982} --- ``Wage Determination and Efficiency in Search Equilibrium''}
%% ------------------------------------------------------------------

\subsection*{Main point}

Wage bargaining and search frictions jointly determine equilibrium unemployment. Two-sided search externalities prevent the decentralised outcome from replicating the social optimum, even when technology is homogeneous.

\subsection*{Model structure}

Steady-state partial equilibrium with fixed-coefficient technology: each filled job produces constant output $y$. Matches form through a stochastic search process; breaks occur at exogenous Poisson rate $b$. Wages are set by symmetric Nash bargaining, splitting the bilateral surplus equally. Asset-value equations for unemployment $W_U$, employment $W_E$, vacant jobs $W_V$, and filled jobs $W_F$ close the model. The equilibrium wage is:
\[
\frac{w}{y} = \frac{rb^{-1} + u^{-1}}{2rb^{-1} + u^{-1} + v^{-1}}.
\]
Search externalities arise because an additional unemployed worker speeds vacancy filling but slows job finding for other workers.

\subsection*{Takeaways for your project}

Diamond (1982) supplies the core conceptual architecture: the bilateral meeting process, asset-value (Bellman) representation of worker and firm states, Nash bargaining wage determination, and the insight that equilibrium wages depend on market tightness $(u,v)$. Your model inherits all of these. Diamond is the building block on which Pissarides (2000) and all subsequent search foundations in this reading list rest.


%% ------------------------------------------------------------------
\section{\citet{Pissarides2000} --- \emph{Equilibrium Unemployment Theory}}
%% ------------------------------------------------------------------

\subsection*{Main point}

The canonical textbook treatment of the DMP model, covering: the matching function, free entry, Nash bargaining, endogenous job destruction, and efficiency. Provides the complete structural skeleton for search-based macro-labour models.

\subsection*{Key chapters relevant to your model}

\begin{description}
    \item[Chapter~1 (DMP core).] Matching function $m(u,v)$ with CRS; tightness $\theta = v/u$; vacancy-filling rate $q(\theta)$; job-finding rate $\theta q(\theta)$. Free-entry job-creation condition: $c/q(\theta) = J^F$. Nash bargaining wages. Beveridge curve: $\delta(1-u) = \theta q(\theta) u$. The Hosios (1990) efficiency condition: $\beta = \eta$ (worker bargaining power equals matching-function unemployment elasticity).

    \item[Chapter~2 (Endogenous job destruction).] Match-specific productivity drawn from a distribution and shocked over time. Matches are destroyed when productivity falls below a \emph{reservation productivity} $R$, determined by $S(R) = 0$. Simultaneous determination of job creation and job destruction. This chapter directly founds your unskilled market's endogenous separation rule $p_U^*$.

    \item[Chapter~4 (On-the-job search).] Workers search on the job and move directly to better matches. Directly relevant to your skilled market.

    \item[Chapter~8 (Efficiency).] With endogenous job destruction, efficiency requires the Hosios condition for creation \emph{and} correct internalisation of the option value of search and the cost of destruction.
\end{description}

\subsection*{Takeaways for your project}

Pissarides (2000) provides the complete structural skeleton of your model: the matching function, free-entry vacancy posting, the Beveridge curve, Nash bargaining, and (from Ch.~2) the endogenous job-destruction mechanism. Market tightness $\theta$ as the key state variable linking job-finding and vacancy-filling rates is adopted throughout your framework.


%% ------------------------------------------------------------------
\section{\citet{postelvinary2002} --- ``Equilibrium Wage Dispersion with Worker and Employer Heterogeneity''}
%% ------------------------------------------------------------------

\subsection*{Main point}

A two-sided heterogeneous search model with on-the-job search in which wages are set through a \emph{sequential auction} (Bertrand competition among employers) generates an endogenous, dispersed cross-sectional wage distribution. This breaks the wage-compression result of Burdett--Mortensen and provides a structural foundation for the empirical wage distribution.

\subsection*{Model structure}

Continuous-time, steady-state. Workers differ in ability $\varepsilon$ (time-invariant); firms differ in productivity $p$; match output is $\varepsilon p$. Unemployed workers receive offers at rate $\lambda_0$; employed workers at rate $\lambda_1$. When an employed worker receives an outside offer, a \emph{poaching competition} (Bertrand auction) begins: the less productive employer reaches its willingness-to-pay ceiling first and the worker either stays at a higher wage or moves. This generates a \emph{two-wage structure} for each worker type: a low entry wage (hired from unemployment) and a high ``promotion'' wage (triggered by outside competition). Estimated on French matched employer--employee data; frictions explain roughly 50\% of wage dispersion.

\subsection*{Takeaways for your project}

PVR is the direct ancestor of your \textbf{skilled market} wage-setting protocol. You inherit: (i)~the sequential auction mechanism, (ii)~the two-wage structure (entry wage when hired from unemployment; renegotiated wage after outside offer), (iii)~the framework for identifying firm and worker heterogeneity distributions, and (iv)~the endogenous wage dispersion driven by the history of outside offers. Your model extends the static PVR setting by adding aggregate shocks and a second (unskilled) market with a separate, simpler Nash bargaining protocol.

\medskip

\noindent\textbf{Important distinction:} your model uses \textbf{Nash bargaining} (CPR 2006 style) in \emph{both} markets, not the full PVR sequential auction. In PVR, firms have all bargaining power over unemployed hires. Your skilled market combines the CPR game-theoretic microfoundation (positive worker bargaining power $\beta > 0$) with the on-the-job Bertrand competition.


%% ------------------------------------------------------------------
\section{\citet{shimer2005} --- ``The Cyclical Behavior of Equilibrium Unemployment and Vacancies''}
%% ------------------------------------------------------------------

\subsection*{Main point}

The standard DMP model with Nash wage bargaining cannot amplify plausible productivity shocks enough to generate observed unemployment and vacancy volatility (the ``Shimer puzzle''). The culprit is fully flexible Nash wages: positive productivity shocks raise workers' outside option, absorbing the gain into wages and leaving vacancy creation nearly unchanged.

\subsection*{Model structure and findings}

Canonical MP model with two aggregate stochastic shocks (labour productivity $p$ and separation rate). Calibrated to U.S. data 1951--2003. Key data facts: the standard deviation of the $v$--$u$ ratio is approximately 20 times that of labour productivity, while the model predicts near-equal volatility. Unemployment and vacancies are strongly negatively correlated ($-0.89$). The job-finding rate co-moves tightly with the $v$--$u$ ratio. The Beveridge curve is the downward-sloping locus of $(u,v)$ pairs consistent with stock-flow balance.

\subsection*{Takeaways for your project}

Shimer (2005) is the benchmark quantitative problem your model is designed to improve upon. You inherit: (i)~the calibration methodology for search models, (ii)~the set of business-cycle moments your model must match (volatility of unemployment, vacancies, $v$--$u$ ratio, job-finding rate), (iii)~the diagnostic that Nash bargaining causes wages to absorb too much productivity variation, and (iv)~the Beveridge curve as an empirical target. Later papers (Robin 2011; Lise and Robin 2017; Moscarini and Postel-Vinay 2024) all take Shimer's puzzle as their point of departure.


%% ------------------------------------------------------------------
\section{\citet{cahuc2006} --- ``Wage Bargaining with On-the-Job Search: Theory and Evidence''}
%% ------------------------------------------------------------------

\subsection*{Main point}

The PVR sequential auction can be given a rigorous game-theoretic foundation via three-player Rubinstein alternating-offers bargaining. This embeds positive worker bargaining power $\beta > 0$ in the wage-setting protocol, making the model empirically estimable with firm-level productivity data.

\subsection*{Model structure}

Continuous time. Workers heterogeneous in ability $\varepsilon$; firms heterogeneous in productivity $p$; match output $\varepsilon p$. Unemployed workers receive $\varepsilon b$ (proportional to ability). Layoff rate $\delta$ exogenous. Key innovation: when an employed worker receives an outside offer, a \emph{three-player Rubinstein alternating-offers game} begins among the worker, current employer, and poaching employer. Under continuous time and equal discount rates, this game has a unique subgame-perfect equilibrium that replicates the PVR sequential auction when $\beta \to 0$ and approaches Nash bargaining when the outside offer is close to the current match value. Estimated on 1993--2000 French matched employer--employee administrative panel; between-firm competition explains more wage gains than Nash bargaining power, especially for unskilled workers (estimated $\beta \approx 0$).

\subsection*{Takeaways for your project}

CPR (2006) is the \textbf{direct structural foundation} for your skilled-market wage equation. You inherit: (i)~the theoretical justification for using the sequential auction when workers have positive $\beta$, (ii)~the identification strategy exploiting both wage and productivity data, (iii)~the result that competition matters quantitatively more than bargaining for low-skill workers (relevant to your unskilled market, where you use standard Nash bargaining), and (iv)~the continuous-time steady-state structure with heterogeneous workers and firms. The CPR wage equation (or your discrete-time analog) is a core building block in your skilled market. Your unskilled market, by contrast, uses the simpler Nash bargaining without the sequential auction.


%% ------------------------------------------------------------------
\section{\citet{robin2011} --- ``On the Dynamics of Unemployment and Wage Distributions''}
%% ------------------------------------------------------------------

\subsection*{Main point}

Embedding PVR-style sequential auctions in a dynamic stochastic setting with aggregate productivity shocks and endogenous job destruction resolves the Shimer puzzle through a \emph{heterogeneity amplification} mechanism: negative shocks push workers at the margin of viable matches into unemployment, generating large, realistic fluctuations from modest productivity volatility.

\subsection*{Model structure}

Discrete time; aggregate state $y_t \in \{y_1 < \cdots < y_N\}$ (finite Markov chain). $M$ worker types with ability $x_m$; firms are \emph{identical} (no firm heterogeneity). Match output: $y_i(m) = x_m y_i$. Both exogenous ($\delta$) and endogenous destruction (when surplus turns negative) operate simultaneously. PVR wage contracting: starting wages at reservation value, renegotiated only when outside offers arrive. Long-term wage contracts create wage stickiness without exogenous rigidity. If 6\% of workers are near the viability margin, unemployment can swing from 4\% to 10\% with only modest productivity volatility. Estimated by simulated GMM. The tails of the wage distribution are more procyclical than middle wages.

\subsection*{Takeaways for your project}

Robin (2011) is your most direct antecedent in the PVR line. You inherit: (i)~the discrete-time, aggregate-shock extension of PVR that founds your aggregate-state Markov chain $z \in \mathcal{Z}$, (ii)~endogenous job destruction as an amplification mechanism (your skilled market's on-the-job search cutoff $p_S^{oj}$ plays an analogous role), (iii)~the two-wage structure under sequential auctions with aggregate shocks, (iv)~the recursive DSGE formulation with a single aggregate state variable, and (v)~the simulated GMM estimation strategy on U.S.\ data. Your paper extends Robin (2011) in three dimensions: adding a second (unskilled) market with a separate equilibrium, endogenising the composition of both labour pools via the training threshold $\bar{x}(z)$, and using Nash bargaining throughout (not the pure PVR monopsony-at-entry protocol).


%% ------------------------------------------------------------------
\section{\citet{bagger2014} --- ``Tenure, Experience, Human Capital, and Wages: A Tractable Equilibrium Search Model of Wage Dynamics''}
%% ------------------------------------------------------------------

\subsection*{Main point}

A tractable equilibrium job search model that jointly estimates human capital accumulation (driven by experience) and employer heterogeneity within the PVR/CPR sequential auction framework. The key tractability device is the piece-rate contract; the key empirical finding is that job shopping (between-job mobility early in the career) is quantitatively more important than human capital accumulation for early wage growth.

\subsection*{Model structure}

Three combined elements: (i)~human capital accumulation (general, driven by experience), (ii)~employer heterogeneity (firms differ in productivity $p$), and (iii)~idiosyncratic productivity shocks. Piece-rate contracts ($w = \phi \cdot \varepsilon_t \cdot p$, piece rate $\phi$ renegotiated via sequential auction upon outside offers) achieve tractability under all three. Estimated on matched employer--employee data using both firm output and worker wages. Wage growth decomposed into: on-the-job human capital accumulation, within-job wage increases from Bertrand competition, and between-job wage increases from job shopping. The ``job-shopping phase'' (first $\sim$10 years) is dominated by between-job mobility.

\subsection*{Takeaways for your project}

If your model features experience-wage profiles or within-market human capital, Bagger et al.\ is the reference. More broadly, you inherit: (i)~the piece-rate device for tractability under idiosyncratic shocks, (ii)~the decomposition of wage growth into human capital and search components, (iii)~the insight that firm heterogeneity and worker heterogeneity together are needed to separately identify rent-sharing from human capital, and (iv)~the estimation methodology using matched data on both output and wages. In your current model, training changes market access but not within-market human capital, so Bagger et al.\ is more of a methodological benchmark than a direct structural predecessor.


%% ------------------------------------------------------------------
\section{\citet{lise2016} --- ``Matching, Sorting and Wages''}
%% ------------------------------------------------------------------

\subsection*{Main point}

An empirical search-matching model with two-sided heterogeneity, productive complementarities, endogenous job destruction, and on-the-job search shows that the AKM firm--worker correlation is a misleading measure of sorting: structural estimation is needed to correctly identify the degree of positive assortative matching. Welfare costs of mismatch are near zero for unskilled workers and substantial for college graduates.

\subsection*{Model structure}

Draws on Mortensen--Pissarides (1994) for endogenous match creation and destruction (separation when idiosyncratic shock drives surplus to zero) and on PVR (2002)/CPR (2006) for wages (Nash bargaining for new matches from unemployment; Bertrand competition for poaching). Key innovation: \emph{productive complementarities} between worker ability $x$ and job type $p$ via a non-separable output function $f(x,p)$, enabling positive assortative matching for output reasons. Estimated on NLSY with simulated method of moments. Welfare decomposition: unskilled labour market near-zero complementarity; college-educated market significant positive complementarity.

\subsection*{Takeaways for your project}

LMR is relevant if your paper deals with sorting or job quality differentials across the two markets. You draw on their framework for: (i)~combining endogenous job destruction (MP) with the sequential auction (PVR/CPR) and on-the-job search, (ii)~the concept of match-specific surplus as the object governing all mobility decisions (block-recursivity), (iii)~welfare decomposition methodology, and (iv)~the treatment of worker--firm complementarities as an identification target. Your model is simpler in that it treats match quality as drawn from a distribution $\Gamma(\cdot \mid z)$ rather than modelling the full $(x,p)$ complementarity function, and adds the education-choice margin that LMR omit.


%% ------------------------------------------------------------------
\section{\citet{lise2017} --- ``The Macro-dynamics of Sorting between Workers and Firms''}
%% ------------------------------------------------------------------

\subsection*{Main point}

Proving block-recursivity in a two-sided heterogeneous search model with aggregate shocks makes tractable the first empirically estimated model of cyclical sorting between workers and firms. Aggregate shocks affect two margins: during recessions, low-type workers matched with high-type firms are fired; during recoveries, low-type workers are re-absorbed.

\subsection*{Model structure}

Extends PVR (2002) with aggregate TFP shocks and extends Robin (2011) with firm heterogeneity and Pissarides-style free-entry vacancy posting. Match output is a function of the worker--firm type pair, allowing complementarities. Firms post state-contingent offers under the sequential auction. Key tractability result: the \emph{match surplus function} does not depend on the current wage, on the distribution of unemployed workers, on the distribution of vacancies, or on the employment distribution --- \textbf{block-recursivity}. The fixed point defining the surplus is solved independently of endogenous distributions. Fitted to aggregate U.S.\ labour market data 1951--2012. The matching set is cone-shaped (strong positive sorting).

\subsection*{Takeaways for your project}

Lise and Robin (2017) is a key technical precedent. You inherit: (i)~the block-recursivity result and the technique of working with the joint surplus rather than separate value functions, (ii)~the method of characterising the matching set (which worker types accept which firm types) under two-sided heterogeneity and aggregate shocks, (iii)~vacancy creation via free entry grafted onto the PVR wage protocol, (iv)~cyclical sorting dynamics as an empirical phenomenon your model should account for, and (v)~the estimation approach combining aggregate time-series data with cross-sectional dispersion moments. Your model takes worker-type distribution as exogenous within each market \emph{conditional} on the training threshold, where Lise and Robin take the full distribution as given. Your paper endogenises the threshold itself.


%% ------------------------------------------------------------------
\section{\citet{acemoglu2001} --- ``Good Jobs versus Bad Jobs''}
%% ------------------------------------------------------------------

\subsection*{Main point}

In a two-type DMP framework (good vs.\ bad jobs differing in sunk creation costs), the decentralised equilibrium produces too few good jobs relative to the social optimum, because good-job firms generate a larger positive externality through rent-sharing that they do not internalise (hold-up problem). Minimum wages and unemployment benefits can improve welfare by correcting job composition.

\subsection*{Model structure}

Extension of the standard DMP model with two job types ($g$ and $b$). Good jobs require larger sunk creation cost $k_g > k_b$; both types sell into competitive markets with imperfect substitution between outputs ($p_g > p_b$). Workers are homogeneous. Search is undirected; matching function $M(u,v)$ with CRS. Wages: asymmetric Nash bargaining with worker power $\beta$. Free entry drives vacancy values to their respective creation costs. The composition of jobs is determined by the relative profitability of posting each type. A hold-up problem arises: good-job firms sink large costs before bargaining, share more rent with workers, and create a larger positive externality that they do not internalise. Multiple equilibria are possible (many good jobs + high outside options; vs.\ many bad jobs + low outside options).

\subsection*{Takeaways for your project}

Acemoglu (2001) is directly relevant to the job-quality dimension of your two-market framework. You inherit: (i)~the two-type job framework within a DMP search environment --- your unskilled market (``bad jobs'') and skilled market (``good jobs'') map directly onto this structure, (ii)~the hold-up logic linking sunk creation costs to equilibrium composition distortions, (iii)~the analysis of how labour market policies affect not just the level but the \emph{composition} of employment, (iv)~the general equilibrium feedback through the worker's outside option, and (v)~the possibility of multiple equilibria. Acemoglu's externality is distinct from Charlot--Decreuse's composition externality but operates through a related channel: in both cases, the marginal entrant into the better market imposes a cost on others that she does not internalise.


%% ------------------------------------------------------------------
\section{\citet{moscarini2024} --- ``On the Job Search and Business Cycles''}
%% ------------------------------------------------------------------

\subsection*{Main point}

Embedding on-the-job search (OJS) with PVR sequential auctions in a DMP business-cycle framework generates three amplification and propagation channels that the standard model entirely lacks: a matching-function elasticity channel, a search-pool composition channel, and an employment-distribution channel. The distribution of employment across the job ladder is an endogenous, slowly-evolving state variable that propagates aggregate shocks.

\subsection*{Model structure}

DMP-style with OJS; both employed and unemployed workers search with different intensities. Firms are ex post heterogeneous in match productivity (drawn upon meeting), generating a job ladder. Free-entry vacancy posting. Wages: PVR sequential auction (Bertrand competition, rank-preserving equilibrium). Because the equilibrium is rank-preserving, the distribution of employment across the job ladder follows a tractable closed-form Markov process. Three OJS channels identified: (1)~\emph{matching-function elasticity} --- accounting for OJS raises estimated vacancy elasticity from 0.32 to 0.50; (2)~\emph{search pool composition} --- employed searchers (more expensive to hire) are more prevalent in booms, dampening job creation; (3)~\emph{employment composition / job ladder} --- after a negative TFP shock, reallocation slows, employment deteriorates (workers clustered in lower-quality matches), gradually stimulating job creation via endogenous persistence.

\subsection*{Takeaways for your project}

MPV (2024) is peripheral but relevant as a high-level orientation. Your model also features OJS in the skilled market and a job ladder generating cyclical employment-distribution dynamics. You inherit from MPV: (i)~the three amplification channels as an organising framework for interpreting your model's cyclical predictions, (ii)~the tractable DMP framework with OJS and free-entry vacancy posting under the sequential auction, (iii)~the rank-preserving equilibrium result, which simplifies computation of the employment distribution, and (iv)~the empirical targets from Shimer (2005) as moments to discipline calibration. If your paper includes a section on business-cycle amplification, MPV provides the direct benchmark. Mark as peripheral in the literature tables.


%% ====================================================================
\newpage
\part{Literature Families}
%% ====================================================================

Below are the papers found, grouped by family, with publication status noted. Papers marked [WP] are working papers and are included for awareness only --- per your instruction, they should not be cited.

\section{Family 1: Education Choice in Search Models}

These are your most direct predecessors. They introduce search frictions into the education decision but do not close both gaps (simultaneous frictional equilibrium in both markets + job ladder in the skilled market + aggregate state transitions).

\begin{longtable}{p{3.8cm} p{3cm} p{3.5cm} p{3.5cm}}
\toprule
\textbf{Paper} & \textbf{Journal} & \textbf{Key idea} & \textbf{Gap your paper fills} \\
\midrule
\endhead
\citet{charlot2005} & \emph{Labour Economics} & Two markets in simultaneous frictional equilibrium, endogenous composition via ability threshold, Nash bargaining, binary education choice & No job ladder in skilled market, no aggregate states, ability-independent schooling cost \\[6pt]
\citet{charlot2005b} & \emph{European Economic Review} & Adaptability and productivity as two margins of education return in a single-technology-segmented market & Single market (no skilled/unskilled segmentation), no binary threshold, no job ladder, no aggregate states \\[6pt]
\citet{wasmer2006} & \emph{AER} & General vs.\ specific human capital in a single frictional market with firing costs; endogenous job destruction & Single market only (no two-market segmentation at all), no composition externality, no job ladder \\[6pt]
\citet{charlot2010} & \emph{Labour Economics} & Two-dimensional heterogeneity (ability + cost), simultaneous two-market equilibrium, over-education for rich / under-education for poor & No job ladder in skilled market, no aggregate states \\[6pt]
\citet{decreuse2013} & \emph{Labour Economics} & Education and hours, joint determination & Same limitations \\[6pt]
\citet{charlot2013} & \emph{Research in Economics} & Education with wage posting and matching [file not yet obtained] & Same limitations \\[6pt]
\citet{agopsowicz2016} & [WP] & Education choice + OJS & Working paper --- do not cite \\
\bottomrule
\end{longtable}

\textbf{Your positioning.}\quad You are the first to combine: (1)~binary education threshold splitting a heterogeneous population, (2)~both markets in simultaneous frictional equilibrium, (3)~a full job ladder (CPR 2006 style) in the skilled market, (4)~endogenous job destruction in the unskilled market, and (5)~aggregate state transitions. Charlot--Decreuse (2005) and (2010) have (1) and (2) but not (3)--(5). Wasmer (2006) has only a single market.


\section{Family 2: Macro / GE Education Choice (Competitive)}

These are the ``big'' education models in competitive frameworks. Your paper positions itself as complementary --- bringing their logic into a frictional world.

\begin{longtable}{p{4cm} p{3.5cm} p{3.5cm} p{3cm}}
\toprule
\textbf{Paper} & \textbf{Journal} & \textbf{Key idea} & \textbf{Gap your paper fills} \\
\midrule
\endhead
\citet{heckman1998} & \emph{Review of Econ.\ Dynamics} & Skill formation, education choice, endogenous skill prices & Competitive markets, no search frictions or unemployment risk \\[6pt]
\citet{abbott2019} & \emph{JPE} & OLG education choice, intergenerational transfers, GE & Same --- competitive factor markets \\[6pt]
\citet{lee2019} & \emph{JPE} & Intergenerational transmission, education + human capital & Same \\[6pt]
\citet{barany2016} & \emph{JoLE} & Education choice + directed technical change + minimum wage & Competitive markets, no search frictions \\
\bottomrule
\end{longtable}


\section{Family 3: Job-Ladder / Search-and-Matching Foundations}

These are the models you build on technically.

\begin{longtable}{p{3.5cm} p{2.8cm} p{4cm} p{3.5cm}}
\toprule
\textbf{Paper} & \textbf{Journal} & \textbf{Key idea} & \textbf{Your extension} \\
\midrule
\endhead
\citet{diamond1982} & \emph{RES} & Search equilibrium, Nash bargaining, bilateral surplus & Foundational building block \\[6pt]
\citet{Pissarides2000} & MIT Press & DMP: matching function, free entry, Beveridge curve, endogenous job destruction & Foundational; you add two-market structure and education margin \\[6pt]
\citet{postelvinary2002} & \emph{Econometrica} & Sequential wage competition, job-ladder model, two-sided heterogeneity & Extended to 2 markets with endogenous composition; Nash bargaining retained (CPR style) \\[6pt]
\citet{shimer2005} & \emph{AER} & Cyclical unemployment dynamics, Shimer puzzle & Motivates aggregate-shock analysis; your model resolves via heterogeneity amplification \\[6pt]
\citet{cahuc2006} & \emph{Econometrica} & Game-theoretic foundation of sequential auction with Nash bargaining power & Direct foundation for your skilled-market wage equation \\[6pt]
\citet{robin2011} & \emph{Econometrica} & PVR with aggregate shocks, endogenous job destruction, amplification via worker heterogeneity & You add second market, education margin, and Nash bargaining throughout \\[6pt]
\citet{bagger2014} & \emph{AER} & Tenure, experience, HC accumulation on the job ladder, piece-rate tractability & HC accumulation within market, not market-selection; methodological benchmark \\[6pt]
\citet{lise2016} & \emph{REDyn} & Two-sided heterogeneity + complementarities + endogenous job destruction + OJS & Adds education margin; your model omits $(x,p)$ complementarity \\[6pt]
\citet{lise2017} & \emph{AER} & Block-recursivity with aggregate shocks and two-sided heterogeneity; cone-shaped matching set & Take type distribution as exogenous within each market; you endogenise it via training threshold \\[6pt]
\citet{acemoglu2001} & \emph{JLE} & Good vs.\ bad jobs, hold-up distortion, policy effects on job composition & You add frictional two-market structure with endogenous education choice \\[6pt]
\citet{moscarini2024} & \emph{Revue écon.} & OJS + business cycles, three amplification channels, employment distribution as state variable & Peripheral; relevant for cyclical dynamics in your skilled market \\
\bottomrule
\end{longtable}


\section{Family 4: Human Capital in Search (Within-Market)}

These study training/HC \emph{within} a given market, not market selection.

\begin{longtable}{p{3.5cm} p{3cm} p{4cm} p{3.5cm}}
\toprule
\textbf{Paper} & \textbf{Journal} & \textbf{Key idea} & \textbf{Distinction from yours} \\
\midrule
\endhead
\citet{bagger2014} & \emph{AER} & HC accumulation on job ladder & Within-market, no education margin \\[6pt]
\citet{audoly2024} & \emph{Fed Staff Reports} & Job-ladder losses, skill depreciation, displaced workers & Same --- within-market \\[6pt]
\citet{flinn2016} & [WP] & Worker-firm joint training decisions & Working paper \\[6pt]
\citet{acabbi2024} & [WP] & Firm-dependent HC, cyclical sorting, hysteresis & Working paper --- interesting but not citable \\
\bottomrule
\end{longtable}


\section{Family 5: Cyclical Enrollment}

Reduced-form evidence that enrollment is countercyclical --- your model provides the structural mechanism.

\begin{longtable}{p{4.5cm} p{4cm} p{5.5cm}}
\toprule
\textbf{Paper} & \textbf{Journal} & \textbf{Key idea} \\
\midrule
\endhead
\citet{dellas2003} & \emph{Oxford Economic Papers} & Countercyclical enrollment, four decades of US data \\[6pt]
\citet{clark2011} & \emph{Economica} & Youth unemployment $\to$ enrollment in England \\[6pt]
\citet{barr2015} & \emph{J.\ Public Economics} & Great Recession $\to$ enrollment at non-selective institutions; UI role \\[6pt]
\citet{bicakova2021} & \emph{Economic Journal} & Graduates who enrol in recessions earn higher wages \\[6pt]
\citet{bicakova2023} & \emph{Labour Economics} & Same authors, extended analysis \\[6pt]
\citet{blom2021} & \emph{JPE} & Community college enrollment responds to local conditions \\[6pt]
\citet{leibing2026} & [WP] & Recent evidence on enrollment cycles + outcomes \\
\bottomrule
\end{longtable}


\section{Family 6: Outside Options and Match Quality}

Empirical and structural work on how outside options shape search outcomes.

\begin{longtable}{p{4.5cm} p{3.5cm} p{6cm}}
\toprule
\textbf{Paper} & \textbf{Journal} & \textbf{Key idea} \\
\midrule
\endhead
\citet{nekoei2017} & \emph{AER} & UI extension improves match quality via selectivity \\[6pt]
\citet{card2018} & \emph{JEEA} & Meta-analysis: training programmes work better medium-run \\[6pt]
\citet{lindner2020} & \emph{AEJ: Applied} & Frontloading vs.\ flow UI --- different search responses \\[6pt]
\citet{cohen2025} & \emph{AEJ: Applied} & Skill depreciation during unemployment \\
\bottomrule
\end{longtable}


%% ====================================================================
\newpage
\part{Papers to Read and Search Strategies}
%% ====================================================================

\section{Papers to Look At More Closely}

Beyond what is already in the literature review, the following papers are worth reading or skimming:

\begin{enumerate}
    \item \textbf{\citet{barr2015}} --- They directly connect UI benefit design to enrollment timing during the Great Recession. Very relevant to your outside-option $\to$ training-threshold story.

    \item \textbf{\citet{bicakova2021}; \citet{bicakova2023}} --- They find that students who enrol during recessions earn more later. Your model could rationalise this: workers who train during recessions are adversely selected on ability (lower $\bar{x}$ $\to$ more marginal workers train) but enter a recovering skilled market.

    \item \textbf{\citet{blom2021}} --- Community college enrollment and local labour markets. Good micro evidence for your countercyclical-enrollment prediction.

    \item \textbf{\citet{cohen2025}} --- Skill depreciation during unemployment. Relevant to why the duration cost of generous outside options matters.

    \item \textbf{\citet{lindner2020}} --- Frontloading vs.\ flow UI. Directly relevant to your Policy~Exercise~(2) (flow training subsidies vs.\ upfront cost reductions).

    \item \textbf{\citet{card2018}} --- Meta-analysis finding that training programmes work better in the medium run. Your model's job-ladder dynamics in the skilled sector provide a structural explanation.
\end{enumerate}


\section{Manual Search Strategies}

\subsection{Strategy 1: Forward Citations from \citet{charlot2005}}

Go to the Google Scholar page for \citet{charlot2005} and click ``Cited by.'' This will surface every paper that builds on the education-in-search tradition. Filter by ``Since 2015'' to find recent work. This is the most targeted way to find papers in your direct predecessor family.

\subsection{Strategy 2: Forward Citations from \citet{postelvinary2002}}

The PVR paper has thousands of citations. Search within them using Google Scholar's search-within-citations feature:
\begin{itemize}
    \item ``Postel-Vinay Robin'' + ``education'' or ``training'' or ``human capital'' or ``skill choice''
    \item This surfaces anyone who extended PVR with an education/training dimension.
\end{itemize}

\subsection{Strategy 3: NBER Working Paper Search}

Go to \url{https://www.nber.org/papers} and search:
\begin{itemize}
    \item ``education choice'' + ``search frictions''
    \item ``job ladder'' + ``education''
    \item ``training'' + ``search model'' + ``equilibrium''
    \item ``outside option'' + ``match quality''
\end{itemize}
NBER working papers that are good enough often end up in top journals. Even if you will not cite the WP itself, it tells you what is in the pipeline and who is working on related questions.

\subsection{Strategy 4: Journal-Specific Searches}

Search within specific top journals for recent papers you might have missed:
\begin{itemize}
    \item \textbf{\emph{Econometrica}}: ``education'' + ``search'' (last 10 years)
    \item \textbf{\emph{AER}}: ``training'' + ``frictional'' or ``search'' (last 10 years)
    \item \textbf{\emph{JPE}}: ``human capital'' + ``job ladder'' (last 10 years)
    \item \textbf{\emph{REStud}}: ``education'' + ``matching'' or ``search frictions'' (last 10 years)
    \item \textbf{\emph{Journal of Labor Economics}}: ``education choice'' + ``unemployment'' (last 10 years)
    \item \textbf{\emph{Labour Economics}}: ``education'' + ``search'' + ``equilibrium'' (last 10 years)
\end{itemize}

\subsection{Strategy 5: Search for ``Education'' in Search-Model Surveys}

Look for recent survey articles on search-and-matching models:
\begin{itemize}
    \item Rogerson, Shimer \& Wright (2005, JEL) --- the classic survey. Check what they say about education/human capital and follow references.
    \item Any recent \emph{Handbook of Labor Economics} chapter on search models.
    \item Wright, Kircher, Julien, Guerrieri (2021, JEL) survey on directed search --- check if education choice appears.
\end{itemize}

\subsection{Strategy 6: Conference Programmes}

Check recent programmes for:
\begin{itemize}
    \item \textbf{SaM (Search and Matching) conference} --- annual conference specifically on search models. Papers presented in 2023--2026 will include unpublished work on the frontier.
    \item \textbf{NBER Summer Institute (Labor Studies)} --- programme available online.
    \item \textbf{CEPR/IZA workshops on Labour Economics} --- programme available online.
    \item \textbf{EEA/ESEM annual meetings} --- search the programme for ``education'' + ``search'' or ``training'' + ``frictional.''
\end{itemize}
Even if the papers are working papers, knowing who is working on what helps you position.

\subsection{Strategy 7: Keyword Combinations on Google Scholar}

Try these specific queries:
\begin{itemize}
    \item \texttt{"endogenous education" "search frictions" "general equilibrium"}
    \item \texttt{"training decision" "job ladder"}
    \item \texttt{"skill choice" "frictional labor market"}
    \item \texttt{"education margin" "unemployment"}
    \item \texttt{"outside option" "reservation wage" "match quality"}
    \item \texttt{"countercyclical enrollment" structural model}
    \item \texttt{"education policy" "search and matching"}
\end{itemize}

\subsection{Strategy 8: Author-Based Search}

These authors work at the intersection of education and search. Check their recent publications and working papers:
\begin{itemize}
    \item \textbf{Olivier Charlot} and \textbf{Bruno Decreuse} --- the education-in-search tradition
    \item \textbf{Fabien Postel-Vinay} --- PVR extensions
    \item \textbf{Jeremy Lise} --- sorting and search
    \item \textbf{Costas Meghir} --- structural labour, education
    \item \textbf{Lance Lochner} --- human capital, education
    \item \textbf{Attila Lindner} --- UI design, outside options (also your advisor's co-author on \citet{lindner2020})
\end{itemize}


%% ====================================================================
\section{Assessment: Where You Stand in the Literature}
%% ====================================================================

\textbf{Unique contribution.}\quad No published paper in a top journal simultaneously:
\begin{enumerate}
    \item Has endogenous education/training choice in a search model, with a binary threshold that splits a heterogeneous population into two distinct labour pools.
    \item Places both markets in frictional equilibrium simultaneously, with separate free-entry conditions, separate tightness levels, and separate composition that feeds back into job creation in each market.
    \item Endogenises the composition of both pools through the training threshold $\bar{x}(z)$.
    \item Builds the skilled market on the CPR (2006) job-ladder framework (Nash bargaining with on-the-job Bertrand competition and on-the-job search cutoff $p_S^{oj}$).
    \item Endogenises job destruction in the unskilled market via a DMP-style reservation-quality rule $p_U^*$.
    \item Introduces aggregate state transitions $z \in \mathcal{Z}$ that shift the full parameter constellation of both markets jointly, making the training threshold $\bar{x}(z)$ cyclically sensitive.
    \item Estimates the model by Simulated Method of Moments on U.S.\ data.
\end{enumerate}

\medskip

\textbf{Positioning relative to direct predecessors.}

\begin{itemize}
    \item \textbf{Charlot--Decreuse (2005)} has (1) and (2) but not (3)--(7). It has endogenous composition via a threshold $\sigma^*$ and two markets in simultaneous equilibrium --- but both markets use static Nash bargaining with no job ladder, no endogenous job destruction (only exogenous death), an ability-independent schooling cost, and no aggregate states. The key distinction is the absence of the job ladder and aggregate dynamics, \emph{not} the absence of simultaneous equilibrium.

    \item \textbf{Charlot--Decreuse-Granier (2005)} has neither (2) nor the binary threshold: it is a representative-agent model in a single (technology-segmented) market. There is no composition externality across two distinct pools.

    \item \textbf{Wasmer (2006)} has a single labour market with no segmentation by skill level. The correct distinction from your paper is \emph{not} ``two markets not simultaneously in equilibrium'' but rather ``single market, no market segmentation at all.'' Wasmer's endogenous job destruction mechanism is useful background for your unskilled market, but the overall architecture is orthogonal.

    \item \textbf{Charlot--Decreuse (2010)} has (1) and (2) via two-dimensional heterogeneity, but again no job ladder, no aggregate states, and symmetric Nash bargaining in both markets. Your paper breaks this symmetry: Nash bargaining in the unskilled market, CPR sequential auction in the skilled market.

    \item \textbf{Charlot--Malherbet (2013):} the correct paper (``Education and Self-Selection in a Search Economy,'' \emph{Research in Economics}) has not yet been obtained. The file in the workspace is a different 3-page note on informality. A summary will be added once the correct paper is located.
\end{itemize}

\medskip

\textbf{Positioning relative to job-ladder foundations.}

\begin{itemize}
    \item The PVR extensions (\citet{lise2017}; \citet{bagger2014}; \citet{robin2011}; \citet{moscarini2024}) have the job-ladder structure but no education margin. Your paper adds the education/training threshold to the job-ladder world --- a dimension none of these papers consider.
    \item The macro education models (\citet{barany2016}; \citet{abbott2019}; \citet{lee2019}) have endogenous education choice and GE feedbacks but in competitive settings. Your paper brings their logic into a frictional world with an explicit job ladder.
\end{itemize}

\medskip

\textbf{Risk areas.}
\begin{itemize}
    \item The \citet{charlot2005} papers are in \emph{Labour Economics} and \emph{European Economic Review}, not in the very top~5. This is actually good for you: it means the question has not been ``solved'' at the top level. But you need to engage carefully with their work and show what is genuinely new --- in particular, that you are not simply adding aggregate shocks to a model that is otherwise isomorphic to Charlot--Decreuse (2005).

    \item The \citet{acabbi2024} WP on cyclical sorting and hysteresis is conceptually close in spirit (aggregate shocks $\to$ composition effects $\to$ hysteresis), but they do not have an education margin. Worth tracking but not currently citable.

    \item Make sure to clearly distinguish from \citet{wasmer2006}: that paper is in a top-5 journal, and referees will ask about it. The key distinction is that Wasmer has a single labour market with no skill-market segmentation --- not ``two markets not simultaneously in equilibrium.''

    \item The file mismatch for Charlot--Malherbet (2013) means that paper's positioning relative to yours cannot yet be assessed. Locate the correct paper before circulating the draft.
\end{itemize}

\bibliographystyle{plainnat}
\bibliography{references}

\end{document}
